\chapter{Various}

\section{Intervals}
	\kactlimport{IntervalContainer.h}
	\kactlimport{IntervalCover.h}
	% \kactlimport{ConstantIntervals.h}
	% \kactlimport{EliminateNestedRanges.h}

\section{Misc. algorithms}
	\kactlimport{TernarySearch.h}
	\kactlimport{LIS.h}
	\kactlimport{FastKnapsack.h}
	% \kactlimport{DecimalBinary.h}
	% \kactlimport{Mex.h}
	% \kactlimport{RandomHashing.h}
	% \kactlimport{Comparators.h}
	\subsection{One-liners}
		\begin{itemize}
			\item Binary: \verb@bitset<64>(x).to_string()@, back: \verb@stoll(s, 0, 2)@.
			\item Set/multiset hashing: give each value a random \verb@rng()@ label, hash = xor (set) or sum (multiset) of labels.
			\item Josephus (every $k$-th dies, 0-indexed survivor): $J(1)=0$, $J(n)=(J(n-1)+k) \bmod n$; for $k=2$: $2(n-2^{\lfloor\log_2 n\rfloor})$.
			\item Inversions: \texttt{oset}, add \verb@i - t.order_of_key(a[i])@ before inserting \texttt{a[i]}.
			\item Nested ranges: sort by $(l\uparrow, r\downarrow)$; range $i$ is contained in an earlier one iff $r_i \le \max r_{<i}$.
		\end{itemize}

\section{Dynamic programming}
	\kactlimport{KnuthDP.h}
	\kactlimport{DivideAndConquerDP.h}

\section{Debugging tricks}
	\begin{itemize}
		\item \verb@signal(SIGSEGV, [](int) { _Exit(0); });@ converts segfaults into Wrong Answers.
			Similarly one can catch SIGABRT (assertion failures) and SIGFPE (zero divisions).
			\verb@_GLIBCXX_DEBUG@ failures generate SIGABRT (or SIGSEGV on gcc 5.4.0 apparently).
		\item \verb@feenableexcept(29);@ kills the program on NaNs (\texttt 1), 0-divs (\texttt 4), infinities (\texttt 8) and denormals (\texttt{16}).
		\item Compile with \verb@-fsanitize=address,undefined -D_GLIBCXX_DEBUG@ to catch out-of-bounds and overflow.
	\end{itemize}

\section{Optimization tricks}
	\verb@__builtin_ia32_ldmxcsr(40896);@ disables denormals (which make floats 20x slower near their minimum value).
	\subsection{Bit hacks}
		\begin{itemize}
			\item \verb@x & -x@ is the least bit in \texttt{x}.
			\item \verb@for (int x = m; x; ) { --x &= m; ... }@ loops over all subset masks of \texttt{m} (except \texttt{m} itself).
			\item \verb@c = x&-x, r = x+c; (((r^x) >> 2)/c) | r@ is the next number after \texttt{x} with the same number of bits set.
			\item \verb@__builtin_popcountll(x)@, \verb@__builtin_ctzll(x)@, \verb@__lg(x)@ (highest set bit).
		\end{itemize}
	\subsection{Pragmas}
		\begin{itemize}
			\item \lstinline{#pragma GCC optimize ("Ofast")} will make GCC auto-vectorize loops and optimizes floating points better.
			\item \lstinline{#pragma GCC target ("avx2")} can double performance of vectorized code, but causes crashes on old machines.
			\item \lstinline{#pragma GCC optimize ("trapv")} kills the program on integer overflows (but is really slow).
		\end{itemize}
	\kactlimport{FastInput.h}
	\kactlimport{BumpAllocator.h}
	% \kactlimport{FastMod.h}
	% \kactlimport{SmallPtr.h}
	% \kactlimport{BumpAllocatorSTL.h}
	% \kactlimport{Unrolling.h}
	% \kactlimport{SIMD.h}
